% \chapter{Overview of Trajectory Prediction Methods and State-of-the-art}
\chapter{Trajectory Prediction: Background and State of the Art}
\label{chapter:sota}

TODO: concrete header for the given thesis's theme!

% Don't write too long about well-known facts in books and text-books.

\section{Taxonomy of Trajectory Prediction Methods}
\label{sec:sota-taxonomy}

In the literature, there are numerous different prediction methods, and several criteria enable their
systematic classification, such as the modality and the representation of the input, the form of the
produced output, or the paradigm by which prediction is integrated with perception. Here the modeling
methods are classified according to the categorization presented in \cite{madjid2025survey}.

\subsection{Physics-Based Methods}
\label{subsec:sota-physics}

These rely on kinematic or dynamic models of motion, using physical equations to extrapolate
future positions and not accouting for uncertainty. Examples include constant-velocity, constant-acceleration models 
(Single-Trajectory simulation) \cite{Lin2000VehicleDynamics} \cite{Miller2002Collision}. 
Kalman filters address the unceratinty limitation of the previous models by modeling the vehicle's state as
a Gaussian disstribution. It predicts trajectories as a two step process of filtering (denoising current
state from noisy measurements) and forecasting \cite{Kalman1960}. Monte Carlo methods, in particular
particle filtering, drop the linearity and Gaussianity assumptions of the Kalman filter by approximating
the state distribution with a set of weighted samples propagated through the motion model, which permits
arbitrarily shaped and multimodal distributions at the cost of a higher computational load
\cite{Gordon1993Bootstrap}. Physics-based approaches are simple and fast,
but they struggle with complex maneuvers or sudden behavior changes since they lack contextual
understanding. Some variants combine multiple motion models to handle different maneuvers
(Multiple-Model methods) \cite{Li2005MultipleModel}.

\subsection{Learning-based methods}
\label{subsec:sota-learning}

Unlike physics-based, learning-based methods learn the motion dynamics from data.
Here, three groups of approaches are mentioned. 

\subsubsection{Classical Machine Learning Methods}
\label{subsec:sota-classical-ml}
These methods learn motion patterns from data using statistical models, but without deep
neural networks. Common approaches include \acp{hmm} \cite{Deo2018Maneuver}, \acp{gp} based prediction
models \cite{Wang2008GPDM}, \acp{dbn} \cite{Li2020Pedestrian},
and \acp{gmm} \cite{Wiest2012GMM}. They can capture
uncertainty to a degree (e.g., an \ac{hmm} can model discrete maneuver states, \acp{gp}
provide a distribution over trajectories) and were popular before deep learning. However, their
capacity to model very complex, high-dimensional scenarios is limited compared to modern
deep nets.

\subsubsection{Reinforcement Learning Methods}
\label{subsec:sota-rl}

\Ac{rl} methods treat prediction as a sequential decision problem. Here, an agent learns a
policy (via reward feedback) to mimic human driving behavior, effectively planning likely
trajectories. Since the reward that drives human driving is not known, \ac{irl}
recovers it from expert demonstrations and uses it to regularise the predictor \cite{Choi2020IRL}, while
\ac{il} skips the reward altogether and extracts the policy directly from the observed
trajectories \cite{Qi2020Imitative}. \Ac{rl}-based predictors can naturally handle multi-agent interactions by
modeling how agents might react to each other, and can ensure physically feasible outputs through learned
policy constraints. However, \ac{rl} methods can be data-hungry, harder to train and stabilize, and
historically have been more common in planning than in pure trajectory forecasting. They are
an active research area for capturing game-theoretic interactions and long-horizon planning in
prediction.

\subsubsection{Deep Learning Methods}
\label{subsec:sota-dl}
Various Deep Learning methods perform successfully in learning motion dynamics and now dominate trajectory
prediction. Their strengths and limitations differ mainly according to their architectural design.



\textbf{Recurrent Neural Networks}

\Acp{rnn} and \ac{lstm} capture temporal history via hidden states updated across timesteps. They are able
to effectively capture temporal patterns for trajectory prediction task. Historically, single recurrent
models were used to map past trajectories to future paths \cite{Hoy2018Pedestrian}, however recent works
facilitate groups or stacks of multiple models \cite{Xin2018DualLSTM} or employ them for one of the subtasks
like detection \cite{Shao2023ReasonNet}. They converge reliably and can be trained autoregressively, in the
same mode in which they are later run at inference, but they lack spatial reasoning and struggle with
long-range temporal dependencies \cite{Liu2021DLSurvey}.

\textbf{Convolutional Neural Networks}

Models built on \acp{cnn} are mostly responsible for feature extraction such as capturing objects' positions
and environmental context \cite{CoverNet2020}. \Acp{tcn} are able to model temporal dependencies through the
causal convolution layers that preserve the temporal order \cite{Nikhil2019CNN}. They extract spatial
features from image-like inputs such as \ac{bev} or rasterised maps efficiently and permit fast inference, 
although the fixed grid on which they operate represents the individual agents only implicitly and is more
expensive than a vectorised encoding of the same scene \cite{Liu2021DLSurvey}.

\textbf{Graph Neural Networks}

\Acp{gnn} represent the scene as a graph whose nodes are the agents and map elements
and whose edges encode the relations between them, which makes them a natural fit for
modelling spatial and temporal interactions \cite{Wang2021GraphTCN}. They excel in dense interaction
scenarios by modelling agent-to-agent relations explicitly, but they add unnecessary complexity in sparse
environments such as highway driving, and their dynamic-graph variants rely on consistent object identities
and therefore inherit the errors of the tracker \cite{madjid2025survey}.

\textbf{Transformers}

Transformers replace recurrence with self-attention, which relates every agent and every map element to
all the others in a single step and thus captures long-range dependencies without having to propagate
information through a hidden state \cite{Liu2021mmTransformer}. Since attending over all agents, time steps
and map elements at once is expensive, practical models factorise the attention along these axes and decode
the future with a set of learned queries that produce several trajectories in parallel
\cite{Wayformer2022}. Their capacity for long-term dependencies makes them particularly promising for
long-horizon prediction, but they are highly sensitive to training stability
and demand extensive tuning and warm-up \cite{madjid2025survey}.
Transformer systems pair a strong scene encoder with multi-modal decoders and explicit
interaction and map context \cite{Liu2021mmTransformer,Ngiam2021SceneTransformer}.

\section{Prediction Paradigms}
\label{sec:sota-paradigms}
In this section, the prediction paradigms are discussed. The prediction paradigm is an important characteristic of a model, as it determines how the
driving stack is decomposed into modules, what is optimised during training, and what
information is passed from perception to prediction. The advantages and the limitations of a given model
largely follow from this choice.


\subsection{Modular (Detect-Track-Predict)}
\label{subsec:sota-modular}
This approach has been of the most common so far implemented in the literature. The paradigm consists 
of sequential steps of detection, tracking and prediction, which were independently under an extensive research. First, the sensor data is captured and
detector localizes the objects present in the scene, usually as a set of bounding boxes with the
corresponding class labels. The tracking module then associates these detections across consecutive frames,
opening a tracklet for every newly observed object and preserving its identity as the scene evolves.
Finally, the prediction module consumes the active tracklets and forecasts the future positions of the corresponding agents \cite{madjid2025survey}.
Even though this sequential decomposition is intuitive and each module can be developed and debugged
separately, the stages are optimised in isolation which causes errors to propagate 
across modules and results in suboptimal performance due to the lack of end-to-end optimization \cite{hu2025survey}.

TODO: add some examples of models

\subsection{Joint Learning (End-to-End Perception \& Prediction)}
\label{subsec:sota-e2e}
Joint learning methods treat detection and motion forecasting as a single optimization problem instead of two independent stages.
These unified optimization goals help to mitigate the error accumulation in the intermediate steps.
Such models are trained with a compound loss over both tasks. They differ in how tightly the stages are coupled,
ranging from a shared \ac{lidar} backbone co-optimized through
a combined loss \cite{Luo2018FaF}, through \ac{bev} features reused by both tasks \cite{Hu2021FIERY}, to full
gradient flow from the prediction loss back through the detector \cite{Liang2020PnPNet}. Query-based models
take this integration further by representing every agent as a learnable query that is iteratively refined
while attending to scene and map features, so that detection, tracking and forecasting are resolved together
\cite{Gu2023ViP3D,Casas2024DeTra}.
The price of this coupling is that task boundaries blur, which makes poor performance hard to attribute to a
specific component, and the iterative refinement of queries adds computational overhead that complicates
real-time deployment in dense scenes \cite{madjid2025survey}.

TODO: add tree of the methods 
% Maybe add if enough space
% \section{Scene Representation and Sensor Fusion}
% \label{sec:sota-fusion}


\section{Datasets}
\label{sec:sota-datasets}

Research in trajectory prediction has been accelerated by large-scale autonomous driving datasets. They
differ in the sensors they record, in the length and the annotation rate of their sequences, and in whether
their annotations support perception, forecasting, or both.

\begin{itemize}
    \item \textbf{KITTI (2012)} \cite{Geiger2012KITTI}: the early \ac{av} benchmark with stereo cameras,
    \ac{lidar} and \ac{gps}/\ac{imu}, widely used for detection, tracking and \ac{slam}. With roughly 150
    short sequences it supports only short-horizon forecasting and contains little social interaction.
    \item \textbf{ApolloScape (2018)} \cite{Huang2018ApolloScape}: Chinese urban driving data covering
    detection, segmentation, tracking and mapping. It is useful for perception and localization, but it is
    less standardized for forecasting than nuScenes or Argoverse.
    \item \textbf{Lyft Level 5 (2019)} \cite{Kesten2019Lyft}: a large urban dataset with synchronized
    \ac{lidar} and cameras and with \acp{hdmap}, similar in format to nuScenes and used mostly for 3D
    detection and tracking.
    \item \textbf{nuScenes (2020)} \cite{Caesar2020nuScenes}: 1{,}000 twenty-second scenes recorded with
    6 cameras, 5 radars and one \ac{lidar} covering $360^\circ$, annotated at \(2\,\mathrm{Hz}\) with 3D
    boxes, \acp{hdmap} and object identities that persist across frames. The official split is 700 training
    and 150 validation scenes, and the prediction challenge adds future trajectories over a
    \(6\,\mathrm{s}\) horizon on exactly these sequences.
    \item \textbf{Argoverse 2 (2021)} \cite{Wilson2021Argoverse2}: the successor of the benchmark that
    popularized map-based prediction and established \acs{ade}, \acs{fde} and \acs{mr} as the reported
    metrics \cite{Chang2019Argoverse}. It is published as several distinct collections, among them a motion
    forecasting set of 250k scenarios with per-scenario \acp{hdmap} and a separate sensor set of 1{,}000
    annotated sequences with seven ring cameras, two stereo cameras and \ac{lidar}. The two collections
    contain different recordings.
    \item \textbf{Waymo Open Motion (2021)} \cite{Ettinger2021WOMD}: more than 100k twenty-second scenarios
    sampled at \(10\,\mathrm{Hz}\), with \acp{hdmap} and a large share of interactive situations. It
    introduced \acs{map} for multi-modal predictions. Later releases of the dataset extend these scenarios
    with \ac{lidar} point clouds \cite{Chen2024WOMDLiDAR} and with camera information, but both cover only
    the first second of each nine-second window, and the camera data is distributed as learned image tokens
    rather than as raw images.
\end{itemize}

NuScenes is used for every experiment reported in this thesis. A jointly trained perception and prediction
model needs synchronized raw \ac{lidar} and camera data, 3D boxes, object identities that persist across
frames so that tracking can be supervised, and future trajectories, all annotated on the same sequences.
Several of the datasets above hold these annotations apart. NuScenes defines an official prediction benchmark on exactly the
sequences it annotates, and it is the dataset on which the end-to-end models discussed in
Section~\ref{sec:sota-paradigms} report their results, which makes the numbers obtained here directly
comparable with them.

TODO: add table with datasets if training on multiple datasets, and also edit metrics if needed


\section{Evaluation Metrics}
\label{sec:sota-metrics}

A jointly trained model has to be judged on every task it solves, so three families of metrics are
reported. The detection metrics measure how well the objects present in the scene are identified, the
tracking metrics measure how reliably their identities are maintained over time, and the trajectory
prediction metrics measure how well their future motion is anticipated. All three families are built on a
notion of a correct match between a prediction and a ground-truth annotation, but they establish that match
in different ways.

\subsection{Detection Metrics}
\label{subsec:sota-metrics-det}

The protocol used throughout this thesis is the one defined with nuScenes \cite{Caesar2020nuScenes} and 
summarized in \cite{Valverde2025Detection}. A
predicted box is matched to a ground-truth box of the same class when the distance between their centres on
the ground plane is below a threshold \(d\), which replaces the intersection over union criterion common in
image-based detection and avoids penalising small objects disproportionately. Matching is greedy over
predictions sorted by decreasing confidence, and every ground-truth box can be matched at most once. The
outcome of this assignment defines the counts:

\begin{itemize}
    \item a \emph{true positive} is a prediction matched to a ground-truth box within threshold \(d\),
    \item a \emph{false positive} is a prediction that finds no such match, either because nothing is there
    or because the object was already claimed by a higher scoring prediction,
    \item a \emph{true negative} is a correctly reported negative, absence of any detection. 
    \item a \emph{false negative} is a ground-truth box left without any matching prediction.
\end{itemize}

Based on these, the following performance metrics are defined:

\begin{enumerate}
    \item \textbf{Precision and recall} express how many of the predictions are correct and how many of the
    annotated objects are found,
    \begin{equation}
        \mathrm{Precision} = \frac{TP}{TP + FP}, \qquad
        \mathrm{Recall} = \frac{TP}{TP + FN}.
        \label{eq:prec-rec}
    \end{equation}
    Sweeping the confidence score threshold traces a precision-recall curve \(p(r)\).

    \item \textbf{\Acf{ap}} summarises that curve for a given class as the area under its interpolated
    form, so that fluctuations of \(p(r)\) at individual operating points do not influence the result
    \cite{Valverde2025Detection},
    \begin{equation}
        \mathrm{AP} = \int_{0}^{1} \max \big\{\, p(r) : r \ge r_{0} \,\big\} \, \mathrm{d}r_{0},
        \label{eq:ap}
    \end{equation}
    where \(p(r)\) is the precision at recall \(r\). The nuScenes protocol evaluates this over the
    operating points that remain once those with a recall or a precision below \(0.1\) have been discarded
    as dominated by noise, and normalises the result to the unit interval \cite{Caesar2020nuScenes}.

    \item \textbf{\Acf{map}} averages \acs{ap} over the set of classes \(\mathcal{C}\) and over a set of
    matching thresholds \(\mathcal{T}\), so that a method is not rewarded for being accurate only at one
    tolerance,
    \begin{equation}
        \mathrm{mAP} = \frac{1}{|\mathcal{C}|\,|\mathcal{T}|}
        \sum_{c \in \mathcal{C}} \sum_{t \in \mathcal{T}} \mathrm{AP}_{c,t},
        \qquad \mathcal{T} = \{0.5,\, 1,\, 2,\, 4\}\,\mathrm{m}.
        \label{eq:map}
    \end{equation}

    \item \textbf{True positive error metrics} describe how good the matched boxes actually are, which
    \acs{map} alone does not capture because a match within \(d\) counts fully regardless of its quality.
    They are computed over the matches obtained at \(d = 2\,\mathrm{m}\) and averaged over classes. They
    comprise the \textbf{translation error ATE} as the centre distance in metres, the \textbf{scale error ASE} as
    \(1 - \mathrm{\acs{iou}}\) after aligning translation and orientation, the \textbf{orientation error AOE} as the
    yaw difference in radians, the \textbf{velocity error AVE} in metres per second, and the \textbf{attribute error AAE} as
    one minus the attribute classification accuracy.

    \item \textbf{\Acf{nds}} combines the two into a single number, weighting the detection quality against
    the five error terms,
    \begin{equation}
        \mathrm{NDS} = \frac{1}{10}\left[\, 5\,\mathrm{mAP} +
        \sum_{\mathrm{mTP} \in \mathbb{TP}} \big(1 - \min(1,\, \mathrm{mTP})\big) \right],
        \label{eq:nds}
    \end{equation}
    where \(\mathbb{TP}\) is the set of the five averaged error metrics. Each error is \textbf{clipped at one}, so an
    arbitrarily bad error cannot drive the score below the contribution of \acs{map} alone.
\end{enumerate}

\subsection{Tracking Metrics}

\label{subsec:sota-metrics-track}

Tracking is as well evaluated with the protocol of the nuScenes tracking benchmark \cite{Caesar2020nuScenes}.
Predictions are matched to ground-truth objects frame by frame with the same centre distance criterion as
in detection, but the assignment now also has to be consistent in time, because a tracker is required to
keep the same identity attached to the same object. The matching therefore proceeds in two stages: a
pairing established in an earlier frame is carried forward whenever both partners are still present and
within the allowed distance, and only the objects and tracks left over are assigned by Hungarian
matching on the centre distance. What is understood as a \emph{true positive} therefore changes:

\begin{itemize}
    \item a \emph{true positive} is a ground-truth object matched within the allowed distance that
    either carried no track before or is matched to the track it already carried,
    \item a \emph{false positive} is a track that is matched to no ground-truth object, so the tracker is
    following something that does not exist,
    \item a \emph{false negative} is a ground-truth object that no track reached in this frame,
    \item a \emph{true negative} is a correctly reported negative, absence of any detection. 
    \item an \emph{\acf{ids}} is a ground-truth object matched to a track whose identity differs from the one it
    carried before. A track that
    is lost and later recovered under a new identity therefore produces an \acs{ids} even though every
    individual box may be accurate.
\end{itemize}

It is also important to state that the lifetime of a trajectory is not simply the set of frames in which it was
annotated. The benchmark first discards annotations that fall outside the
per-class evaluation range or that contain no \acs{lidar} returns at all, and
then linearly interpolates each trajectory across any remaining interior gap,
so that a fully occluded object is restored. A trajectory's lifetime is therefore the uninterrupted span between
its first and last surviving observation, and a tracker is expected to carry
an identity through occlusions that the sensor itself did not observe. The tracking metrics are as follows:

\begin{enumerate}
    \item \textbf{\Acf{mota}} aggregates all the ways a tracker can fail into a single number, counting
    every false positive, every missed object and every identity switch against the number of annotated
    object instances \(GT\) summed over all frames,
    \begin{equation}
        \mathrm{MOTA} = 1 - \frac{FP + FN + \mathrm{IDS}}{GT}.
        \label{eq:mota}
    \end{equation}
    The measure is not bounded from below, as a tracker that reports far more objects than exist can drive
    it to an arbitrarily negative value.

    \item \textbf{\Acf{motp}} describes how precisely the matched objects are localised, as the mean centre
    distance \(d_i\) over the set of true positives,
    \begin{equation}
        \mathrm{MOTP} = \frac{1}{|TP|} \sum_{i \in TP} d_i .
        \label{eq:motp}
    \end{equation}

    \item \textbf{\Acf{amota} and \acf{amotp}} remove the dependence on the confidence threshold in the same
    way \acs{ap} does for detection. The tracking score is swept, and the two measures above are averaged
    over a set \(\mathcal{R}\) of forty recall levels, from which those below \(0.1\) are discarded as
    dominated by noise,
    \begin{equation}
        \mathrm{AMOTA} = \frac{1}{|\mathcal{R}|} \sum_{r \in \mathcal{R}} \mathrm{MOTAR}_r ,
        \label{eq:amota}
    \end{equation}
    where \(\mathrm{MOTAR}_r\) is the value of \acs{mota} at recall \(r\), normalised so that it occupies
    the unit interval at every recall level and is therefore comparable across them. \Acs{amota} is the
    primary measure of the benchmark.

    \item \textbf{Supporting counts} describe the failure modes that the aggregate scores hide. Unlike
    \acs{mota}, whose unit is a single object instance in a single frame, these are evaluated along the
    whole lifetime of one annotated object and each trajectory is counted once. 
    A trajectory is reported as \textbf{mostly tracked (MT)} when it is followed for at least \(80\%\) of its
    lifetime and as \textbf{mostly lost (ML)} when it is followed for at most \(20\%\). The remaining
    trajectories are \textbf{partially tracked (PT)}, so that the three categories partition the annotated
    trajectories,
    \begin{equation}
        N_{\mathrm{traj}} = \mathrm{MT} + \mathrm{PT} + \mathrm{ML},
        \label{eq:traj-partition}
    \end{equation}
    where each count refers to whole trajectories and \(N_{\mathrm{traj}}\) are all annotated trajectories that survive the evaluation filters.

     A \textbf{fragmentation} is counted 
    whenever a correctly followed trajectory is interrupted and later resumed. The
    \textbf{track initialisation duration} expresses how long it takes before an object is picked up for
    the first time, and the \textbf{longest gap duration} how long it remains lost in the worst
    interruption.
\end{enumerate}

These supporting counts matter for a model whose tracks are gated by a confidence threshold. A detector
whose individual boxes are accurate but whose scores fluctuate around that threshold loses and reacquires
its objects, which leaves \acs{ap} almost unaffected while raising the number of fragmentations and
identity switches and therefore lowering \acs{amota}.


\subsection{Trajectory Prediction Metrics}
\label{subsec:sota-metrics-pred}
The prediction metrics are computed over agents. An agent is characterised only by its position in the
ground plane and its forecast trajectory. The predicted agents for a frame are assigned to the \ac{gt}
agents of that frame by Hungarian matching on the centre distance, with the same allowed distance of
\(2\,\mathrm{m}\) as in Section~\ref{subsec:sota-metrics-track}.
The track identity is never taken into account, so the assignment is recomputed from scratch in
every frame rather than carrying established pairings as in tracking. The confidence of a prediction is likewise not read, so the counts are
evaluated at the single operating point at which the tracker emitted its tracks.

A matched agent is additionally only useful if its forecast is close to what actually happened, which is
what the following metrics quantify. For an agent with \(K\) predicted future trajectories
\(\hat{y}^{(k)}\), a ground-truth future \(y\) and the set \(\mathcal{T}\) of future time steps with a valid
annotation:

\begin{enumerate}
    \item \textbf{min\acs{ade}} is the average displacement of the best of the \(K\) modes over the whole
    horizon, which rewards a model for containing one good hypothesis without punishing it for the
    diversity of the others,
    \begin{equation}
        \mathrm{minADE}_K = \min_{k \in \{1,\dots,K\}}
        \frac{1}{|\mathcal{T}|} \sum_{t \in \mathcal{T}}
        \big\lVert \hat{y}^{(k)}_t - y_t \big\rVert_2 .
        \label{eq:minade}
    \end{equation}
    A single annotated future step is sufficient for an agent to contribute, and the average is taken over
    \(|\mathcal{T}|\) rather than over the full horizon, so agents that leave the scene early are scored on
    a correspondingly shorter forecast.

    \item \textbf{min\acs{fde}} evaluates only the endpoint of the horizon, where the uncertainty of the
    forecast is largest,
    \begin{equation}
        \mathrm{minFDE}_K = \min_{k \in \{1,\dots,K\}}
        \big\lVert \hat{y}^{(k)}_{T} - y_{T} \big\rVert_2 .
        \label{eq:minfde}
    \end{equation}
    It is defined only for agents that are still annotated at the end of the horizon, that is
    \(T \in \mathcal{T}\). An agent that disappears earlier is excluded entirely 

    \item \textbf{\Acf{mr}} turns the endpoint error into a binary outcome and reports how often none of the
    \(K\) modes lands near the true endpoint,
    \begin{equation}
        \mathrm{MR} = \frac{\big|\{\, i : \mathrm{minFDE}_K^{(i)} > \tau \,\}\big|}{N_{\mathrm{surv}}},
        \qquad \tau = 2\,\mathrm{m},
        \label{eq:mr}
    \end{equation}
    where the index \(i\) runs over the \(N_{\mathrm{surv}}\) matched agents that survive the whole horizon,
    since the others have no min\acs{fde} to threshold.

    \item \textbf{\Acf{epa}} \cite{Gu2023ViP3D} is the metric that judges the two tasks jointly. An agent
    contributes to the numerator only if it was both matched and forecast accurately, and every false
    positive is subtracted with a penalty \(\alpha\),
    \begin{equation}
        \mathrm{EPA} = \max\left(0,\,
        \frac{N_{\mathrm{hit}} - \alpha\, N_{FP}}{N_{\mathrm{GT}}}\right),
        \qquad \alpha = 0.5,
        \label{eq:epa}
    \end{equation}
    where \(N_{\mathrm{hit}}\) counts the matched agents whose \(\mathrm{minFDE}_K\) does not exceed
    \(\tau\), and \(N_{\mathrm{GT}}\) the annotated agents that survive the whole horizon. Unlike the three
    metrics above, which are averaged over the agents that were found, \acs{epa} is normalised by the
    annotated agents rather than the matched ones, so \emph{a model cannot improve it by detecting fewer objects
    and forecasting only the easy ones}.
\end{enumerate}

The metrics are complementary and have to be read together. A model may reach a low min\acs{ade} while
missing many agents entirely, because the error is averaged only over the matched ones, and \acs{mr}
exposes the cases where no mode comes close at all. \Acs{epa} penalises both failure modes at once and is
therefore used as the primary criterion in this work. 